\begin{lemma}[Differentiating Mehler's formula: $\partial_j P_s f = e^{-s}P_s(\partial_j f)$ on the smooth core]
Fix $s > 0$ and write $r = e^{-s} \in (0,1)$ and $\tau = \sqrt{1-e^{-2s}} \in (0,1)$.
For a function $g : \mathbb{R}^n \to \mathbb{R}$ put
\[
  (A_s g)(x) = \int_{\mathbb{R}^n} g\bigl(rx + \tau z\bigr) \, d\gamma_n(z),
\]
the Mehler average of \noderef{MehlerFormula}, with $\gamma_n$ as in
\noderef{GaussMeasure}.  Let $f : \mathbb{R}^n \to \mathbb{R}$ be smooth with
compact support and let $j \in \{1,\ldots,n\}$.  Then:
\begin{enumerate}
\item for every $x$, the function $z \mapsto f(rx+\tau z)$ is
$\gamma_n$-integrable;
\item for every $x$, the function $z \mapsto (\partial_j f)(rx+\tau z)$ is
$\gamma_n$-integrable;
\item the function $A_s(\partial_j f)$ is continuous on $\mathbb{R}^n$; and
\item for every $x$ the function $A_s f$ is differentiable at $x$, and
\[
  \partial_j (A_s f)(x) = e^{-s}\,\bigl(A_s(\partial_j f)\bigr)(x) .
\]
\end{enumerate}

Clause (4) is the reference paper's one-clause remark that, for
$f \in C_c^\infty$, differentiation of Mehler's formula gives
$DP_sf = e^{-s}P_s Df$.  The statement is in the Euclidean normalization
$\partial_j$ rather than in the energy normalization $D_j = 2^{-1/2}\partial_j$,
and the two are the same identity: multiplying both sides of (4) by the constant
$2^{-1/2}$ turns $\partial_j$ into $D_j$ on the left and $\partial_j f$ into
$D_j f$ on the right, with no change of constant.  The connection to the
semigroup itself is \noderef{MehlerFormula}: for every $v \in L^2(\gamma_n)$ the
class $P_s v$ agrees $\gamma_n$-almost everywhere with $A_s$ applied to a
representative of $v$, so $A_s f$ is a representative of $P_s f$ for the class of
$f$ and $A_s(\partial_j f)$ one of $P_s(\partial_j f)$.

Clauses (1)--(3) are stated rather than left implicit because every later use
needs them: (1) and (2) make the two averages genuine integrals rather than the
zero value the Bochner convention assigns to a non-integrable integrand, and
(3), combined with (4) applied at each coordinate, is what makes $A_s f$
continuously differentiable and therefore available for ordinary integration by
parts.
\end{lemma}
\begin{proof}
Fix $s>0$, so $0 < r = e^{-s} < 1$, hence $0 < 1-e^{-2s} < 1$ and
$0 < \tau < 1$; in particular $r > 0$ and $\tau > 0$.

\emph{Step 0: $f$ and its first derivatives are bounded.}
The function $f$ is continuous with compact support, so it is bounded: writing
$S = \operatorname{supp} f$, which is compact, $|f|$ attains a maximum on $S$ by
the extreme value theorem (shared context) and vanishes off $S$; put
$M_0 = \sup_{\mathbb{R}^n}|f| < \infty$.  Since $f$ is smooth, its total
derivative $Df : \mathbb{R}^n \to (\mathbb{R}^n)^*$ is continuous, and it
vanishes off $S$: if $x \notin S$ then $f$ vanishes identically on an open
neighbourhood of $x$, so all its derivatives vanish at $x$.  Hence
$\operatorname{supp}(Df) \subseteq S$ is compact and, by the same extreme value
argument applied to the continuous function $y \mapsto \|Df(y)\|$,
$M = \sup_{\mathbb{R}^n}\|Df\| < \infty$.  In particular
$M_1 := \sup_{\mathbb{R}^n}|\partial_j f| \leq M < \infty$, and $Df$, being
continuous with compact support, is uniformly continuous on $\mathbb{R}^n$
(shared context: a continuous function with compact support is uniformly
continuous, by Heine--Cantor on a compact neighbourhood of its support together
with vanishing outside the support).

\emph{Step 1: clauses (1) and (2).}
Fix $x$.  The map $z \mapsto rx+\tau z$ is continuous, so $z \mapsto f(rx+\tau z)$
is continuous, hence measurable, and it is bounded by $M_0$.  By
\noderef{GaussIsProb} the measure $\gamma_n$ is a probability measure, hence
finite, and a bounded measurable function on a finite measure space is
integrable.  This is clause (1).  The same argument with $\partial_j f$, which is
continuous and bounded by $M_1$ by Step 0, gives clause (2).

\emph{Step 2: clause (3).}
Let $x_k \to x$ in $\mathbb{R}^n$.  For every fixed $z$ we have
$rx_k + \tau z \to rx+\tau z$, so $(\partial_j f)(rx_k+\tau z) \to
(\partial_j f)(rx+\tau z)$ by continuity of $\partial_j f$.  All these functions
of $z$ are bounded in absolute value by the constant $M_1$, which is
$\gamma_n$-integrable because $\gamma_n$ is finite (\noderef{GaussIsProb}).
Dominated convergence (shared context) gives
$(A_s(\partial_j f))(x_k) \to (A_s(\partial_j f))(x)$.  As $\mathbb{R}^n$ is a
metric space, sequential continuity is continuity, which is clause (3).

\emph{Step 3: differentiation under the integral sign.}
Fix $x$ and define the linear map
\[
  \Lambda_x : \mathbb{R}^n \to \mathbb{R},
  \qquad
  \Lambda_x(h) = r\int_{\mathbb{R}^n} Df(rx+\tau z)[h] \, d\gamma_n(z) .
\]
This is well defined and bounded: for each $h$ the integrand is continuous in
$z$ and bounded by $M\,|h|$, hence $\gamma_n$-integrable as in Step 1, and
linearity in $h$ follows from linearity of $Df(y)$ at each $y$ together with
linearity of the integral; moreover $|\Lambda_x(h)| \leq rM|h|$, so $\Lambda_x$
is a bounded linear functional.

We claim $A_s f$ is differentiable at $x$ with derivative $\Lambda_x$.  Let
$\varepsilon > 0$.  By the uniform continuity of $Df$ (Step 0) choose
$\delta > 0$ with $\|Df(a) - Df(b)\| \leq \varepsilon$ whenever $|a-b| \leq \delta$.
Let $h \in \mathbb{R}^n$ with $0 < |h| \leq \delta/r$ (recall $r > 0$).  For each
fixed $z$, apply the fundamental theorem of calculus to the smooth function
$t \mapsto f(rx+\tau z + t\,rh)$ on $[0,1]$:
\[
  f\bigl(r(x+h)+\tau z\bigr) - f(rx+\tau z)
    = \int_0^1 Df(rx+\tau z + t\,rh)[rh] \, dt ,
\]
using $r(x+h)+\tau z = (rx+\tau z) + rh$.  Subtracting
$Df(rx+\tau z)[rh] = \int_0^1 Df(rx+\tau z)[rh]\,dt$,
\[
  \Bigl| f\bigl(r(x+h)+\tau z\bigr) - f(rx+\tau z) - Df(rx+\tau z)[rh] \Bigr|
    \leq \int_0^1 \bigl\|Df(rx+\tau z+t\,rh) - Df(rx+\tau z)\bigr\|\,|rh|\,dt
    \leq \varepsilon\, r|h| ,
\]
because $|t\,rh| \leq r|h| \leq \delta$ for $t \in [0,1]$.  This bound is uniform
in $z$.  Integrating in $z$ against the probability measure $\gamma_n$ --- each
of the three terms being $\gamma_n$-integrable, the first two by Step 1 and the
third by the first paragraph of this step --- and using that $|\int\cdot| \leq
\int|\cdot|$ together with $\gamma_n(\mathbb{R}^n) = 1$,
\[
  \bigl|(A_s f)(x+h) - (A_s f)(x) - \Lambda_x(h)\bigr| \leq \varepsilon\,r\,|h| .
\]
Since $\varepsilon > 0$ was arbitrary, the difference quotient error is
$o(|h|)$, i.e.\ $A_s f$ is differentiable at $x$ with derivative $\Lambda_x$.

\emph{Step 4: clause (4).}
Evaluating the derivative at the $j$th standard basis vector $e_j$ and using
$Df(y)[e_j] = (\partial_j f)(y)$,
\[
  \partial_j(A_s f)(x) = \Lambda_x(e_j)
    = r\int_{\mathbb{R}^n} (\partial_j f)(rx+\tau z)\,d\gamma_n(z)
    = e^{-s}\,\bigl(A_s(\partial_j f)\bigr)(x) ,
\]
which is clause (4); the factor $r = e^{-s}$ is exactly the inner derivative of
the affine map $x \mapsto rx + \tau z$, which is where the paper's $e^{-s}$ comes
from.
\end{proof}
